\chapter{Complex Laguerre geometry and transverse convexity}\label{ch:transverse-convexity}

\section{Classical theorem route}

The XF-4 scalar admits a second theorem-level interpretation that is independent of the Dimitrov--Xu density route.  Csordas and Varga developed necessary and sufficient Laguerre conditions for the Riemann Xi function, and Csordas and Escassut state the complex Laguerre criterion in a form that makes its geometry explicit \cite{CsordasVarga1990,CsordasEscassut2005}.

Let
\[
\Xi(z)=\xi\!\left(\frac12+iz\right),
\qquad z=x+iy.
\]
The completed Xi function is real entire of order one, and its zeros lie in the horizontal strip $|\Im z|<1/2$ before any RH specialization.  Thus Xi belongs to the strip class required by the complex Laguerre criterion.

\section{Transverse curvature identity}

Define
\begin{equation}\label{eq:xf5-modulus-surface}
M(x,y)=|\Xi(x+iy)|^2.
\end{equation}
For a real entire function,
\[
\overline{\Xi(x+iy)}=\Xi(x-iy).
\]
Using
\[
\partial_y\Xi(x+iy)=i\Xi'(x+iy),
\qquad
\partial_y^2\Xi(x+iy)=-\Xi''(x+iy),
\]
direct differentiation of Eq.~\eqref{eq:xf5-modulus-surface} gives
\begin{align}
\partial_y^2M(x,y)
={}&2|\Xi'(x+iy)|^2\\
&-2\Ree\!\left[
\Xi(x+iy)\overline{\Xi''(x+iy)}
\right].
\end{align}
With the XF-4 scalar $Q_\Xi$ from Eq.~\eqref{eq:q-xi-definition},
\begin{equation}\label{eq:xf5-curvature-identity}
\boxed{
\partial_y^2|\Xi(x+iy)|^2=2Q_\Xi(x,y).
}
\end{equation}
The repository records Eq.~\eqref{eq:xf5-curvature-identity} as \status{EXACT\_PASS}.

\section{Complex Laguerre criterion}

For a real entire function $f$ in the strip class $S(A)$, the classical complex Laguerre criterion states
\begin{equation}\label{eq:complex-laguerre-general}
f\in\mathcal{LP}
\iff
|f'(z)|^2
\ge
\Ree\!\left[f(z)\overline{f''(z)}\right]
\qquad\forall z\in\mathbb C.
\end{equation}
Csordas and Escassut explicitly identify the geometric content of Eq.~\eqref{eq:complex-laguerre-general}: for every fixed real $x$, the function $|f(x+iy)|^2$ is convex in $y$ \cite{CsordasEscassut2005}.

Applying Eq.~\eqref{eq:complex-laguerre-general} to Xi and using Eq.~\eqref{eq:xf5-curvature-identity} yields
\begin{equation}\label{eq:xi-lp-convexity}
\boxed{
\Xi\in\mathcal{LP}
\iff
\partial_y^2|\Xi(x+iy)|^2\ge0
\quad\forall x,y\in\mathbb R.
}
\end{equation}

\section{RH equivalence}

For the Riemann Xi function, membership in the Laguerre--Pólya class is equivalent to all Xi zeros being real, hence to RH.  Therefore Eq.~\eqref{eq:xi-lp-convexity} gives the classical geometric criterion
\begin{theorem}[Transverse-convexity criterion]\label{thm:xf5-transverse-convexity}
Under the standard real-entire strip-class hypotheses for Xi,
\begin{equation}\label{eq:xf5-rh-convexity}
\boxed{
\mathrm{RH}
\iff
\forall x,y\in\mathbb R:\quad
\partial_y^2|\Xi(x+iy)|^2\ge0.
}
\end{equation}
Equivalently,
\begin{equation}\label{eq:xf5-rh-q-nonnegative}
\boxed{
\mathrm{RH}
\iff
Q_\Xi(x,y)\ge0
\quad\forall x,y\in\mathbb R.
}
\end{equation}
\end{theorem}

The equivalence in Theorem~\ref{thm:xf5-transverse-convexity} is classified \status{STANDARD}.  The universally quantified curvature condition is classified
\[
\status{OPEN\_RH\_EQUIVALENT\_CRITERION}.
\]

\section{Relation to the Wiener--Laguerre route}

Theorem~\ref{thm:wiener-laguerre-rh} gives, through Dimitrov--Xu and Wiener,
\[
\mathrm{RH}
\iff
Q_\Xi(x,y)>0
\quad
\forall x\in\mathbb R,
\quad0<|y|<\frac12.
\]
The complex Laguerre theorem gives instead
\[
\mathrm{RH}
\iff
Q_\Xi(x,y)\ge0
\quad
\forall x,y\in\mathbb R.
\]
The distinction at $y=0$ is essential: the real-axis Laguerre expression may vanish at a multiple real zero, whereas the XF-4 criterion is posed away from the real axis and uses strict positivity.

Thus XF-4 and XF-5 provide two independently sourced theorem routes to the same scalar.  The repository preserves both provenance chains in the typed solver.

\section{Executable check}

The function
\begin{center}
\texttt{xi\_transverse\_modulus\_curvature(x, y)}
\end{center}
returns the analytic expression $2Q_\Xi(x,y)$.  The regression suite independently differentiates $|\Xi(x+iy)|^2$ twice with respect to $y$ at reference points and checks agreement with Eq.~\eqref{eq:xf5-curvature-identity}.

Finite evaluations retain \status{NUMERICAL\_DIAGNOSTIC} authority.  The global criterion in Eq.~\eqref{eq:xf5-rh-convexity} remains the analytic frontier.

\section{Next structural target}

The new representation changes the form of the open obligation from Fourier nonvanishing to curvature positivity:
\begin{equation}\label{eq:xf5-frontier}
\boxed{
\partial_y^2|\Xi(x+iy)|^2\ge0
\quad
\text{for every }(x,y)\in\mathbb R^2.
}
\end{equation}
A useful next step must therefore derive this sign from an independently controlled structure, for example a positive kernel representation, total positivity, or a positive operator form whose quadratic form equals the transverse curvature.
